\section{Pilot Results: Historical Validation}
\label{sec:results}

We ran the full protocol on the ten frozen
\texttt{evidence\_graph\_v1} questions. Headline numbers: every
question was substantively engaged by the 2021--2026 literature
(coverage 100\%); two were answered, seven partially addressed, one
independently posed and still open; one premise was refuted. Evidence
strength: 2.4 supporting papers per question on average, with 60\% of
labels supported by $\geq 2$ independent papers. The earliest
judge-cited engagement came 1--5 years after the cutoff (mean 2.9).
Table~\ref{tab:perq} gives the per-question picture.

\begin{table}[t]
  \centering
  \small
  \caption{Per-question outcomes for \texttt{evidence\_graph\_v1} on
  Astronomy v1 (system rank order; $|E|$ = independent supporting
  papers; year = earliest cited engagement).}
  \label{tab:perq}
  \begin{tabular}{@{}clllcc@{}}
    \toprule
    Rank & ID & Outcome & Premise & $|E|$ & Year \\
    \midrule
    1  & q\_002 & \labelfmt{posed\_but\_open}      & \labelfmt{still\_plausible} & 1 & 2022 \\
    2  & q\_001 & \labelfmt{partially\_addressed}  & \labelfmt{supported}        & 3 & 2021 \\
    3  & q\_004 & \labelfmt{partially\_addressed}  & \labelfmt{supported}        & 1 & 2021 \\
    4  & q\_008 & \labelfmt{answered}              & \labelfmt{refuted}          & 3 & 2025 \\
    5  & q\_007 & \labelfmt{partially\_addressed}  & \labelfmt{still\_plausible} & 3 & 2024 \\
    6  & q\_010 & \labelfmt{partially\_addressed}  & \labelfmt{still\_plausible} & 1 & 2022 \\
    7  & q\_005 & \labelfmt{answered}              & \labelfmt{supported}        & 2 & 2024 \\
    8  & q\_009 & \labelfmt{partially\_addressed}  & \labelfmt{supported}        & 5 & 2021 \\
    9  & q\_006 & \labelfmt{partially\_addressed}  & \labelfmt{still\_plausible} & 1 & 2025 \\
    10 & q\_011 & \labelfmt{partially\_addressed}  & \labelfmt{still\_plausible} & 4 & 2024 \\
    \bottomrule
  \end{tabular}
\end{table}

\subsection{Case study: a premise refuted (q\_008, HD 209458 b)}
\label{sec:results:case}

From pre-2021 evidence, the system flagged a single-dataset conclusion:
the influential retrieval of a \emph{strongly subsolar} water abundance
at the terminator of HD~209458\,b \citep{macdonald2017}. The frozen
question asked, in 2020 terms:

\begin{quote}\itshape
Is the strongly subsolar terminator water abundance retrieved for
HD~209458\,b a property of the atmosphere or an artifact of retrieval
assumptions, as tested by comparing independent retrieval frameworks on
the same and on independent datasets?
\end{quote}

The 2021--2026 literature then performed exactly the test the question
specified (Figure~\ref{fig:case}): a reanalysis of HST and JWST spectra
with improved systematics treatment and Bayesian model averaging, an
independent retrieval framework demonstrating that free-vs-equilibrium
chemistry assumptions span the subsolar-to-solar range, and ground-based
high-resolution spectroscopy constraining the abundance independently of
space-based data. All three converge: the terminator water abundance is
consistent with solar, and the strongly subsolar value was an artifact
of earlier retrieval assumptions and data systematics. Under the
taxonomy this is \labelfmt{answered} + \labelfmt{refuted}: the question
was resolved \emph{by} the community overturning the premise the
question challenged---an outcome that no expert score assigned in 2020
could have certified, and precisely what backtesting exists to detect.

\begin{figure}[t]
\centering
\begin{tikzpicture}[
    font=\small,
    evt/.style={align=center, font=\footnotesize},
    arr/.style={-{Stealth[length=2.2mm]}, thick}]

  \draw[arr] (0,0) -- (13.2,0);
  \node[font=\footnotesize] at (13.0,-0.35) {year};
  \foreach \x/\y in {0.8/2017, 3.0/2020, 5.2/2021, 7.4/2023, 9.6/2025, 11.8/2026}
    \draw (\x,0.08) -- (\x,-0.08) node[below, font=\footnotesize] {\y};

  \node[evt, text width=2.6cm] at (0.8,1.35)
    {subsolar H$_2$O\\ retrieved for\\ HD 209458 b};
  \draw[arr] (0.8,0.75) -- (0.8,0.12);

  \draw[fill=blue!15, rounded corners=2pt] (2.2,1.6) rectangle (3.8,2.35);
  \node[evt] at (3.0,1.975) {frozen\\ question};
  \draw[arr] (3.0,1.6) -- (3.0,0.12);
  \node[evt] at (3.0,2.65) {\emph{artifact or atmosphere?}};

  \draw[dashed] (3.0,-0.7) -- (3.0,1.4);
  \node[evt] at (3.0,-1.15) {cutoff\\ 2020-12-31};

  \node[evt, text width=3.4cm] at (9.6,1.45)
    {three independent\\ reanalyses converge:\\ H$_2$O consistent with solar};
  \draw[arr] (9.6,0.75) -- (9.6,0.12);
  \node[evt, text width=3.9cm] at (9.6,-0.85)
    {\labelfmt{answered} + premise \labelfmt{refuted}};

  \draw[thick, orange!80!black] (5.2,-1.7) -- (11.8,-1.7)
    node[midway, below, font=\footnotesize\itshape]
    {future corpus window (2021--2026)};
  \draw[thick, orange!80!black] (5.2,-1.78) -- (5.2,-1.62);
  \draw[thick, orange!80!black] (11.8,-1.78) -- (11.8,-1.62);
\end{tikzpicture}
\caption{Timeline of q\_008. The question was frozen from pre-2021
evidence; by 2025 three independent analyses had performed the test it
specified and refuted its premise---the strongly subsolar water
abundance was a retrieval artifact.}
\label{fig:case}
\end{figure}

\subsection{Secondary observations}

\paragraph{The top-ranked question is independently posed and open.}
The system's rank-1 question (q\_002: how terminator heterogeneity
biases the WASP-12\,b water abundance and C/O ratio) tracks a
methodological concern the community has since engaged in general form
--- inhomogeneous-terminator biases in retrievals --- without resolving
it for WASP-12\,b specifically: \labelfmt{posed\_but\_open}. A question
the field poses but has not answered is a live research target; that the
system's top pick lands there is the behavior a ranking is supposed to
produce, though $n=1$ at rank 1 proves nothing by itself.

\paragraph{An answered null result.} q\_005 asked whether HST
transmission data independently support NH$_3$ or HCN in HD~209458\,b
--- a molecular-detection claim from the same single-dataset analysis as
q\_008 \citep{macdonald2017}. By 2024--2025, high-resolution
spectroscopy had placed stringent upper limits on both species:
\labelfmt{answered}, premise \labelfmt{supported} (the data indeed do
not independently support the detection). Backtesting counts a
cleanly resolved null exactly as it counts a positive.

\paragraph{Instrument-gated engagement.} q\_011 (whether JWST validated
pre-launch predictions of TRAPPIST-1 CO$_2$ detectability,
\citealp{lustigyaeger2019}) could not have been engaged before JWST flew;
its first cited engagement is 2024 and stellar contamination has so far
prevented a definitive test. Engagement timing is partly an instrument
schedule, not purely a question-quality signal---a confound
Section~\ref{sec:threats} treats explicitly.

\paragraph{What these ten questions do not show.} With ten questions
from one system in one domain, rates carry wide intervals (the exact
95\% Clopper--Pearson interval for 10/10 coverage is $[0.69, 1.0]$),
the baseline rows are judge-only rather than adjudicated
(Section~\ref{sec:baselines}), and the generating system's LLM
components postdate the cutoff (Section~\ref{sec:threats}). The pilot
demonstrates that the protocol runs end-to-end, yields auditable,
reproducible labels, and prices its baselines; it does not establish
that any system reliably anticipates future science.
